\lecture{9}
\title{Reducibility}

\begin{document}

\subsection{Reducibility}

\begin{theorem}{Halting Problem}
    The language $\mathrm{HALT}_{\mathrm{TM}} :=
    \{\langle M, w \rangle \mid M \text{ is a TM that halts on } w\}$
    is undecidable.
\end{theorem}

\begin{proof}
    Our strategy will be to use the undecidability of
    $A_{\mathrm{TM}} := \{\langle M, w \rangle \mid
    M \text{ is a TM that accepts } w\}$ from
    \href{/18.404/lectures/8/#a-tm-undecidability}{Lecture 8}.

    Suppose FTSOC that there is a TM $R$
    that decides $\mathrm{HALT}_{\mathrm{TM}}$.
    Then one can construct $S$
    that does the following on $\langle M, w \rangle$.
    \begin{itemize}
        \item Run $R$ on $\langle M, w \rangle$,
        which, in finite time, must determine whether
        $M$ halts on $w$.
        \begin{itemize}
            \item If $M$ is known to halt on $w$, then
            $S$ should just run $M$ on $w$ and return the result.
            \item If $M$ is known not to halt on $w$, then
            $S$ should just return \textsc{Reject}.
        \end{itemize}
    \end{itemize}
    So the decidability of $\mathrm{HALT}_{\mathrm{TM}}$
    implies the decidability of $A_{\mathrm{TM}}$,
    contradiction. ~ $\blacksquare$
\end{proof}

The above proof is an example
of a broader proof technique: \textbf{reduction.}

\textbf{Definition.}
We say that a language $A$ is \textbf{reducible} to a language $B$
if we can decide $A$ using a decider for $B$.

\begin{theorem}{Reducibility and Decidability}
    Suppose language $A$ is reducible to language $B$.
    Then the following are true:
    \begin{itemize}
        \item If $B$ is decidable, then $A$ is decidable.
        \item If $A$ is undecidable, then $B$ is undecidable.
    \end{itemize}
\end{theorem}

None of this is particularly nonobvious.

\begin{theorem}{$E_{\mathrm{TM}}$}
    The language $E_{\mathrm{TM}} := \{ \langle M \rangle \mid
    M \text{ is a TM and }  L(M) = \emptyset \}$
    is undecidable.
\end{theorem}

\begin{proof}
    Recall from \href{/18.404/lectures/8/#a-tm-undecidability}{Lecture 8}
    that $A_{\mathrm{TM}}$ is undecidable.
    It therefore suffices to prove that
    $A_{\mathrm{TM}}$ reduces to $E_{\mathrm{TM}}$.

    Suppose we have a decider $R$ for $E_{\mathrm{TM}}$.
    Then we can construct a decider $S$ for $A_{\mathrm{TM}}$
    that does the following on $\langle M, w \rangle$.
    \begin{itemize}
        \item Construct a Turing machine $M'$ that
        does the following on input $w'$.
        \begin{itemize}
            \item If $w' \neq w$, instantly \textsc{Reject}.
            \item If $w' = w$, run the instructions of $M$
            (on input $w' = w$).
        \end{itemize}
        \item Now $S$ just returns the opposite of the result of $R$
        on input $\langle M' \rangle$.
        ~ $\blacksquare$
    \end{itemize}
\end{proof}

\subsection{Mapping Reducibility}

\textbf{Definition.}
We say that a language $A$ is \textbf{mapping reducible} to a language $B$
(denoted $A \leq_m B$) if there is a computable function
(i.e., computable via Turing machine)
$f \colon \Sigma^* \to \Sigma^*$ such that
$w \in A$ if and only if $f(w) \in B$ for all $w \in \Sigma^*$.

\begin{theorem}{Mapping Reducibility and Turing-Recognizability}
    If $A \leq_m B$, then the following are true.
    \begin{itemize}
        \item If $B$ is Turing-recognizable,
        then $A$ is Turing-recognizable.
        \item If $A$ is Turing-unrecognizable,
        then $B$ is Turing-unrecognizable.
        \item If $B$ is decidable, then $A$ is decidable.
    \end{itemize}
\end{theorem}

Again, none of this is particularly nonobvious.

\begin{theorem}{$EQ_{\mathrm{TM}}$}
    The language $EQ_{\mathrm{TM}} := \{\langle B, C \rangle \mid
    B \text{ and } C \text{ are TMs with } L(B) = L(C)\}$
    is not Turing-recognizable.
\end{theorem}

\begin{proof}
    Recall from \href{/18.404/lectures/8/#turing-unrecognizable}{Lecture 8}
    that $\overline{A_{\mathrm{TM}}} := \{\langle M, w \rangle \mid
    M \text{ is a TM that does not accept } w\}$ is Turing-unrecognizable.

    It therefore suffices to prove that
    $\overline{A_{\mathrm{TM}}} \leq_m EQ_{\mathrm{TM}}$.
    In other words, we seek a map $f$ such that:
    \[
    f \colon \langle M, w \rangle \mapsto \langle B, C \rangle
    ~~~ \text{ such that } ~~~
    \boxed{\langle M, w \rangle \in \overline{A_{\mathrm{TM}}}
    \iff \langle B, C \rangle \in EQ_{\mathrm{TM}}}
    \]
    We'll define the image $\langle B, C \rangle$
    of $f$ on $\langle M, w \rangle$ like so.
    \begin{itemize}
        \item Take $B$ to always return \textsc{Reject}.
        \item Take $C$ to be the Turing machine that
        just runs $M$ on $w$ (independent of the input $x$).
        \begin{itemize}
            \item Note that $C$ might \textsc{Accept},
            or \textsc{Reject}, or even loop.
        \end{itemize}
    \end{itemize}
    You can check for yourself that
    the boxed equivalence ($\iff$) holds for this $f$.
    ~ $\blacksquare$
\end{proof}

And one more example.

\begin{theorem}{$\overline{EQ_{\mathrm{TM}}}$}
    The language $\overline{EQ_{\mathrm{TM}}} := \{\langle B, C \rangle \mid
    B \text{ and } C \text{ are TMs with } L(B) \neq L(C)\}$
    is not Turing-recognizable.
\end{theorem}

\begin{proof}
    Again, it suffices to prove that
    $\overline{A_{\mathrm{TM}}} \leq_m \overline{EQ_{\mathrm{TM}}}$.
    In other words, we seek a map $f$ such that:
    \[
    f \colon \langle M, w \rangle \mapsto \langle B, C \rangle
    ~~~ \text{ such that } ~~~
    \boxed{\langle M, w \rangle \in \overline{A_{\mathrm{TM}}}
    \iff \langle B, C \rangle \in \overline{EQ_{\mathrm{TM}}}}
    \]
    The image $\langle B, C \rangle$ of $\langle M, w \rangle$ under $f$
    is almost entirely the same as before.
    \begin{itemize}
        \item Take $B$ to always return \textsc{Accept}.
        \item Take $C$ to be the Turing machine that
        just runs $M$ on $w$ (independent of the input $x$).
        \begin{itemize}
            \item Note that $C$ might \textsc{Accept},
            or \textsc{Reject}, or even loop.
        \end{itemize}
    \end{itemize}
    The only thing that changed is that
    $B$ now always returns \textsc{Accept} instead of \textsc{Reject}.
    Check for yourself that this works. ~ $\blacksquare$
\end{proof}

\end{document}